\documentclass{article}
\usepackage[english]{babel}

\title{ForTheL Demo (LaTeX Mode)}
\author{Test File}
\date{2026}

\newenvironment{forthel}{}{}
\newenvironment{signature}{\noindent\textbf{Signature.}}{}
\newenvironment{definition}{\noindent\textbf{Definition.}}{}
\newenvironment{axiom}{\noindent\textbf{Axiom.}}{}
\newenvironment{theorem}[1][]{\noindent\textbf{Theorem}\ifthenelse{\equal{#1}{}}{}{~(#1)}\textbf{.}}{}
\newenvironment{lemma}[1][]{\noindent\textbf{Lemma}\ifthenelse{\equal{#1}{}}{}{~(#1)}\textbf{.}}{}

\begin{document}
\maketitle

This document demonstrates ForTheL in LaTeX mode.

\begin{forthel}
  \begin{signature}
    A group is an object.
  \end{signature}

  \begin{signature}
    Let $G$ be a group.
    An element of $G$ is an object.
  \end{signature}

  \begin{signature}
    Let $G$ be a group and $a, b$ be elements of $G$.
    $a \cdot b$ is an element of $G$.
  \end{signature}

  \begin{signature}
    Let $G$ be a group.
    $e_G$ is an element of $G$.
  \end{signature}

  \begin{axiom}
    Let $G$ be a group and $a$ be an element of $G$.
    Then $e_G \cdot a = a$ and $a \cdot e_G = a$.
  \end{axiom}

  \begin{signature}
    Let $G$ be a group and $a$ be an element of $G$.
    $a^{-1}$ is an element of $G$.
  \end{signature}

  \begin{axiom}
    Let $G$ be a group and $a$ be an element of $G$.
    Then $a \cdot a^{-1} = e_G$ and $a^{-1} \cdot a = e_G$.
  \end{axiom}

  \begin{axiom}
    Let $G$ be a group and $a, b, c$ be elements of $G$.
    Then $(a \cdot b) \cdot c = a \cdot (b \cdot c)$.
  \end{axiom}
\end{forthel}

With the group axioms established, we can prove basic properties.

\begin{forthel}
  \begin{lemma}[Uniqueness of Identity]
    Let $G$ be a group and $a, e$ be elements of $G$.
    If $a \cdot e = a$ then $e = e_G$.
  \end{lemma}
  \begin{proof}
    Assume $a \cdot e = a$.
    Then $e = e_G \cdot e = (a^{-1} \cdot a) \cdot e = a^{-1} \cdot (a \cdot e)
         = a^{-1} \cdot a = e_G$.
  \end{proof}

  \begin{theorem}[Inverse is Involutive]
    Let $G$ be a group and $a$ be an element of $G$.
    Then $(a^{-1})^{-1} = a$.
  \end{theorem}
  \begin{proof}
    $(a^{-1})^{-1}
      = (a^{-1})^{-1} \cdot e_G
      = (a^{-1})^{-1} \cdot (a^{-1} \cdot a)
      = ((a^{-1})^{-1} \cdot a^{-1}) \cdot a
      = e_G \cdot a
      = a$.
  \end{proof}
\end{forthel}

\end{document}
